%======================================================================
%  Companion to Chapter 2 -- Elliptic problems: diffusion and elasticity
%======================================================================
\chapter{Elliptic problems: diffusion and elasticity}\label{cc:ch2}

\framepart{Outline}

Chapter~\ref{B-ch:elliptic} of the notes ends with the plate with a hole,
computed with linear triangles in Python and compared with Kirsch's
solution. Its part~(d) asks for the same computation with Cast3M and other
elements. This chapter does it with four element types and three meshes, and
adds the check that the support reactions balance the applied load. You
should then be able to run a convergence study inside one gibiane file, with
the mesh built in a procedure.

\section{Where Cast3M enters Chapter 2}\label{cc:sec2-map}
\index{Cast3M!Chapter 2}%

\begin{gbox}{Chapter 2 of the notes and Cast3M}\label{cc:tab2}
\small
\begin{tabular}{@{}p{42mm}@{\hspace{3mm}}p{98mm}@{}}
Kirsch, Ex.~\ref{B-exo:kirsch} (d)
 & \textbf{Exercise \ref{exo:cc-kirsch}}: \texttt{TRI3}, \texttt{TRI6}, \texttt{QUA4},
   \texttt{QUA8}; $n_\theta=12,24,48$; $\sigma_{11}(0,a)$, $\sigma_{22}(a,0)$.\\[1mm]
Multipliers as reactions, Sec.~\ref{B-ssec:ritz-lagrange},
Ex.~\ref{B-exo:lagrange-el}
 & \textbf{Exercise \ref{exo:cc-kirsch}}: the reactions on $x=0$ sum to $-sL$.\\[1mm]
Diffusion Lagrangian, Sec.~\ref{B-ssec:diff-lagrangian}
 & To do: the thermal model with an imposed temperature; \op{reac} gives the
   flux $\mu=\vect{p}\cdot\vn$ through the boundary, and its sum balances the
   source.\\[1mm]
Zero-energy modes, Ex.~\ref{B-exo:q-modes}
 & To do: one free \texttt{QUA4} element, the eigenvalues of its stiffness
   (\op{vibr} with a unit mass): three zero modes.\\[1mm]
Anisotropy, Ex.~\ref{B-exo:stiffness}
 & To do: \op{mate} \texttt{ANISOTROPE} with the constants of cadmium, a
   homogeneous test giving $s_{13}$ and $s_{33}$.\\
\end{tabular}
\end{gbox}

\section{A convergence study in one file}\label{cc:sec2-loop}
\index{convergence!Cast3M}%
\begin{castemcom}
The mesh of the quarter plate is built in a procedure whose arguments are the
element type and the number of cells along the hole. \op{opti elem} inside
the procedure sets the type of the elements that \op{droi}, \op{cerc} and
\op{dall} create: \texttt{SEG2} lines give \texttt{TRI3} or \texttt{QUA4}
surfaces, \texttt{SEG3} lines \texttt{TRI6} or \texttt{QUA8}.

The two loops then rebuild the model, the stiffness and the solution for each
mesh. \op{resu} sums the reactions of the symmetry line $x=0$: whatever the
mesh, they balance the traction $s$ on the edge $x=L$, $\sum F_x=-sL$.
\end{castemcom}
\begin{castemlines}
\begin{verbatim}
debproc plaque typ*mot nt*entier;
  opti elem typ;
  ...
finproc dom l12 l23 l45 p1 p5;

repeter bty (dime ltyp);
  typ = extr ltyp &bty;
  repeter bnt (dime lnt);
    nt = extr lnt &bnt;
    dom l12 l23 l45 p1 p5 =
        plaque typ nt;
    ...
    u  = reso (kk et bx et by) ff;
    rx = resu (exco (reac bx u) 'FX');
  fin bnt;
fin bty;
\end{verbatim}
\end{castemlines}

\framepart{Summary}

\begin{gbox*}
\begin{itemize}
\item \textbf{A mesh built in a procedure} turns a convergence study into two
  loops: element type, mesh size.
\item \textbf{The element type is a mesh option}: \op{opti elem} decides
  linear or quadratic, triangles or quadrangles.
\item \textbf{The reactions check the computation}: their sum balances the
  applied load on every mesh, an exact identity of the discrete problem
  (Exercise~\ref{B-exo:lagrange-el}).
\end{itemize}
\end{gbox*}

\framepart{Exercises}

\exercise{Kirsch with four elements}\label{exo:cc-kirsch}
\index{Kirsch problem!Cast3M}%
Take the quarter plate of Exercise~\ref{B-exo:kirsch} ($L=20a$, plane
stress, $E=1$, $\nu=0.3$, tension $s=1$ on $x=L$). For \texttt{TRI3},
\texttt{TRI6}, \texttt{QUA4} and \texttt{QUA8} and $n_\theta=12$, 24, 48 cells
on the quarter hole, compute the nodal stresses $\sigma_{11}(0,a)$ and
$\sigma_{22}(a,0)$, the number of unknowns, and the sum of the reactions on
$x=0$. How many unknowns does each element need for $1\%$ on the stress
concentration factor?
\begin{solution}
Reference values with linear triangles (Python,
\texttt{kirsch\_fem.py}): $\sigma_{11}(0,a)=2.970$, 3.013, 3.020 and
$\sigma_{22}(a,0)=-0.722$, $-0.871$, $-0.946$ for $n_\theta=12$, 24, 48,
against Kirsch's 3 and $-1$ (the plate is finite: with $L=160a$ the factor
is $3.00$). The reactions sum to $-sL=-20$ on every mesh. Expect the
quadratic elements to reach $1\%$ on $\sigma_{11}(0,a)$ with far fewer
unknowns than \texttt{TRI3}; the run gives the numbers. The compressive
$\sigma_{22}(a,0)$ converges slowly with linear triangles (still $5\%$ off at
$n_\theta=48$): it is the harder test.
\untested{\texttt{cast3m/ch2/kirsch\_elements.dgibi}}
\lstinputlisting[style=gibstyle,firstline=22]{cast3m/ch2/kirsch_elements.dgibi}
\end{solution}

\begin{chapterbib}{9}
\bibitem{cc-kirsch} G.~Kirsch, Die Theorie der Elastizit\"at und die
  Bed\"urfnisse der Festigkeitslehre, \emph{Zeitschrift des Vereines deutscher
  Ingenieure} 42 (1898) 797--807.
\end{chapterbib}
